% ============================================================
% Chapter 3 - Methodology (rewritten to match Chapter 4)
% Needs in your preamble:  \usepackage{amsmath,booktabs}
% Every \fill{...} is something only YOU know: replace it with the
% real detail, or delete the sentence if it is not true for your study.
% ============================================================
\providecommand{\fill}[1]{\textbf{[#1]}}

\chapter{Methodology}

This chapter presents the methodology used to develop and evaluate a predictive
framework for informal settlement expansion in Windhoek. The study follows a
data-driven approach that integrates spatial analysis and machine learning. The
chapter describes the research design, the study area and data, preprocessing,
feature engineering, model specification, training and validation, evaluation
metrics, spatial risk mapping, the software used, and ethical considerations.

\section{Research Design}

The study adopts a quantitative, predictive research design grounded in data
science principles. The study area is divided into a regular grid of spatial
cells, and the task is framed as a binary classification problem: for each cell,
estimate the probability that informal settlement expansion occurs. Two
supervised learning models are developed and compared on identical data and
metrics. Logistic regression serves as an interpretable baseline that quantifies
the direction and strength of each spatial driver, and a Random Forest serves as
the main predictive model, able to capture nonlinear effects and interactions. A
single decision tree is included as a benchmark. The fitted probabilities are
then projected across the grid to produce a spatial vulnerability map.

\section{Study Area and Data}

\subsection{Study Area and Spatial Unit}

The study area is the Windhoek metropolitan region, bounded by latitudes
$22.65^{\circ}$S to $22.45^{\circ}$S and longitudes $17.00^{\circ}$E to
$17.18^{\circ}$E. The area is divided into a regular grid of $50 \times 50 =
2{,}500$ cells. Each cell is one observation and represents a localised land
parcel evaluated for informal settlement expansion.

\subsection{Data Sources}

The predictors and the response variable were obtained from publicly available
spatial and socio-economic sources, summarised in Table~\ref{tab:vars}.

\begin{table}[ht]
\centering
\caption{Variables used in the study}
\label{tab:vars}
\begin{tabular}{llll}
\toprule
Variable & Description & Unit & Source / derivation \\
\midrule
RoadDist    & Distance to nearest major road          & km        & \fill{source} \\
SettleDist  & Distance to nearest informal settlement & km        & \fill{source} \\
PopDensity  & Population density                      & persons/km$^2$ & \fill{source} \\
Elevation   & Terrain elevation                       & m (ASL)   & \fill{DEM source} \\
Slope       & Terrain slope                           & degrees   & derived from \fill{DEM} \\
CBDDist     & Distance to the Central Business District & km      & \fill{reference point} \\
InfraProx   & Distance to nearest essential infrastructure & km   & \fill{what counts as infrastructure} \\
Expansion   & Informal settlement expansion (response) & 0/1      & \fill{how labels were derived} \\
\bottomrule
\end{tabular}
\end{table}

\fill{State how the response variable was derived, for example change detection
between two named dates using satellite imagery or settlement footprints, and
state the time period it covers. If the dataset was simulated or generated
rather than observed, say so clearly here.}

\section{Data Preprocessing}

Spatial data were converted into a tabular format in which each grid cell is a
row. Missing values were handled by \fill{method, or state that none were
present}. The response variable $Y_i$ is binary:
\begin{equation}
Y_i =
\begin{cases}
1 & \text{if informal settlement expansion was observed in cell } i,\\
0 & \text{otherwise.}
\end{cases}
\end{equation}
Of the 2{,}500 cells, 416 (16.6\%) are expansion cells, so the classes are
imbalanced; stratified sampling is therefore used so that this proportion is
preserved in every training, test and validation subset.

Predictors were standardised using the mean and standard deviation of the
training set only, to avoid information leaking from the test set:
\begin{equation}
z_{ij} = \frac{x_{ij} - \bar{x}_j}{s_j},
\end{equation}
where $\bar{x}_j$ and $s_j$ are the training mean and standard deviation of
predictor $j$.

\section{Feature Engineering}

Seven spatial predictors were derived for every grid cell, selected for their
relevance to urban expansion in the literature: distance to major roads
(\texttt{RoadDist}), distance to existing informal settlements
(\texttt{SettleDist}), population density (\texttt{PopDensity}), elevation
(\texttt{Elevation}), slope (\texttt{Slope}), distance to the Central Business
District (\texttt{CBDDist}) and distance to essential infrastructure
(\texttt{InfraProx}). Distances were computed from cell centroids using
\fill{GIS or Python method}, and slope was derived from the digital elevation
model. The predictors were combined with the response variable into a single
dataset. Pearson correlation coefficients and group-wise boxplots were used
in the exploratory analysis to examine relationships among predictors and with
the response.

\section{Model Specification}

\subsection{Logistic Regression}

Let $p_i = P(Y_i = 1)$ denote the probability of expansion in cell $i$. The
logistic regression model is
\begin{equation}
p_i = \frac{1}{1 + e^{-z_i}}, \label{eq:logistic}
\end{equation}
with linear predictor
\begin{equation}
z_i = \beta_0 + \beta_1\,\text{RoadDist}_i + \beta_2\,\text{SettleDist}_i
+ \beta_3\,\text{PopDensity}_i + \beta_4\,\text{Elevation}_i
+ \beta_5\,\text{Slope}_i + \beta_6\,\text{CBDDist}_i
+ \beta_7\,\text{InfraProx}_i, \label{eq:linpred}
\end{equation}
which can be written in logit form as
\begin{equation}
\log\!\left(\frac{p_i}{1-p_i}\right) = \beta_0 + \sum_{j=1}^{7}\beta_j X_{ji}.
\label{eq:logit}
\end{equation}
Because the predictors are standardised, $e^{\beta_j}$ is the odds ratio
associated with a one-standard-deviation increase in $X_j$, holding the other
predictors constant. Coefficients are estimated by maximum likelihood, and
standard errors, $z$-statistics and $p$-values are used to assess significance
at the 5\% level.

\subsection{Random Forest}

The Random Forest is an ensemble of $B$ decision trees trained on bootstrap
samples, with a random subset of predictors considered at each split. The
class prediction is the majority vote and the estimated probability of
expansion is the share of trees voting for expansion:
\begin{equation}
\hat{Y}_i = \text{majority vote}\{T_1(X_i), T_2(X_i), \ldots, T_B(X_i)\},
\qquad
\hat{p}_i = \frac{1}{B}\sum_{b=1}^{B} \mathbb{I}\big(T_b(X_i) = 1\big),
\label{eq:rf}
\end{equation}
where $T_b(X_i)$ is the prediction of the $b$-th tree and $\mathbb{I}(\cdot)$ is
the indicator function. This study uses $B = 100$ trees with a maximum depth of
8. Variable importance is measured by the mean decrease in Gini impurity,
\begin{equation}
i(t) = 1 - \sum_{k} p_k(t)^2, \qquad
\text{Imp}(X_j) = \frac{1}{B}\sum_{b=1}^{B}\sum_{t \in T_b:\, v(t)=j}
\frac{N_t}{N}\,\Delta i(t),
\label{eq:gini}
\end{equation}
where $p_k(t)$ is the proportion of class $k$ at node $t$, $v(t)$ is the
variable used to split node $t$, $N_t$ is the number of observations reaching
the node and $N$ the total number of observations.

\subsection{Benchmark Model}

A single decision tree trained on the same data and predictors is used as a
benchmark, to show how much the ensemble improves on one tree.

\section{Model Training and Validation}

The 2{,}500 cells were split into a training set of 2{,}000 cells (80\%) and a
holdout test set of 500 cells (20\%) using stratified random sampling and a
fixed random seed for reproducibility. Model parameters were estimated on the
training set, and predictive performance was assessed on the test set. To check
robustness, 5-fold stratified cross-validation was also applied to the full
dataset, and the mean and standard deviation of the ROC-AUC across folds are
reported.

Because neighbouring grid cells are spatially autocorrelated, a random split can
give optimistic performance estimates. \fill{Optional: include only if you run
it. As a sensitivity check, spatial block cross-validation was also performed,
in which contiguous blocks of cells are held out together.}

\section{Model Evaluation}

Predicted classes are obtained with a probability threshold of 0.5. Performance
is summarised with a confusion matrix (true positives TP, false positives FP,
true negatives TN and false negatives FN) and the following metrics:
\begin{align}
\text{Accuracy}  &= \frac{TP + TN}{TP + TN + FP + FN}, &
\text{Precision} &= \frac{TP}{TP + FP}, \\
\text{Recall}    &= \frac{TP}{TP + FN}, &
F_1 &= \frac{2\,\text{Precision}\times\text{Recall}}{\text{Precision}+\text{Recall}}.
\end{align}
The Receiver Operating Characteristic (ROC) curve plots the true positive rate
$TP/(TP+FN)$ against the false positive rate $FP/(FP+TN)$ over all thresholds,
and the Area Under the Curve (ROC-AUC) summarises how well a model separates the
two classes. Precision-recall curves are also reported because expansion cells
are the minority class.

\section{Spatial Prediction and Risk Mapping}

The best-performing model is applied to every grid cell to obtain an estimated
expansion probability $\hat{p}_i$, which is plotted at the cell coordinates to
produce a continuous vulnerability surface for Windhoek. For interpretation,
cells are grouped into five risk classes, as shown in Table~\ref{tab:classes}.

\begin{table}[ht]
\centering
\caption{Risk classes used for spatial vulnerability mapping}
\label{tab:classes}
\begin{tabular}{ll}
\toprule
Risk class & Probability range \\
\midrule
Critical high-risk         & $\hat{p} > 0.85$ \\
High risk                  & $0.70 < \hat{p} \le 0.85$ \\
Moderate expansion buffer  & $0.30 \le \hat{p} \le 0.70$ \\
Low risk                   & $0.10 \le \hat{p} < 0.30$ \\
Very low risk / protected  & $\hat{p} < 0.10$ \\
\bottomrule
\end{tabular}
\end{table}

\section{Tools and Software}

The analysis was implemented in Python. Pandas and NumPy were used for data
manipulation, Scikit-learn for the machine learning models and evaluation,
Statsmodels for the logistic regression coefficient estimates, standard errors
and $p$-values, and Matplotlib and Seaborn for visualisation.
\fill{Name the GIS software or libraries actually used for spatial processing,
for example QGIS, GeoPandas or Google Earth Engine.}

\section{Ethical Considerations}

The study uses publicly available data and no sensitive personal information.
Predictions are produced for grid cells, not for individuals or households, and
are intended to support planning rather than enforcement against residents. The
risk maps are decision-support tools that carry uncertainty, so results are
reported with appropriate caveats to avoid misrepresentation or misuse.

\section{Chapter Summary}

This chapter described the grid-based predictive framework, the seven spatial
predictors, the logistic regression, Random Forest and decision tree models,
the training and validation design, the evaluation metrics, and the risk
mapping procedure. The next chapter presents the results obtained from
implementing this methodology.
